\subsection*{一步一步建立直观}

\paragraph{1. 把信息量看成排除可能性的成本}

若两种结果等可能，一个二选一问题就能区分它们。若其中一种很罕见，常见结果可以在平均意义上花更少的描述成本，罕见结果则需要更长描述。右图从公平二元分布出发，基准熵为一比特；对数中的底二对应二进制问题的单位。



\paragraph{2. 平均代价要用真实发生概率加权}

把真实分布改为不均匀。罕见结果的单次信息量变大，但它出现得少，所以熵要把每种信息量乘以真实概率再相加。不能只看最大码长就判断平均压缩难度。这里显示的是理想平均编码代价，有限样本的整数码长还受具体编码方案影响。



\paragraph{3. 分布的不确定性与模型失配不同}

保持真实分布固定，改用偏斜的模型。蓝色曲线升高，橙色熵基线不变，说明额外成本来自模型失配，而不是数据本身突然更随机。熵与交叉熵应分开理解；前者由真实分布决定，后者还依赖你拿来编码或预测的模型。
